% ============================================================================
% Chapter 8 --- Network and segmentation
% ----------------------------------------------------------------------------
% Purpose : zones, segmentation strategy, attachment service, Science DMZ.
%           Belongs to the Execution plane (D5.1). Largest layer chapter
%           because of the zone catalogue and the SASE discussion.
% Page target: ~7 pages.
% Dependencies: Ch. 5 (standards: RADIUS, NETCONF/YANG, BGP/EVPN, eBPF/CNI,
%               IPFIX), Ch. 7 (policy decisions are enforced here for east-west
%               and attachment), Ch. 3 (sovereignty framework).
% Mirrors layer chapter template established in Ch. 6 (D6.5).
% ----------------------------------------------------------------------------
% Decisions registered in _coherence-EN.md:
%   D8.1 minimum compliance set for the network layer (the contract)
%   D8.2 reference instantiation = VLAN/VRF + FreeRADIUS + OvS + Science DMZ
%   D8.3 alternatives = proprietary macro-segmentation, micro-segmentation overlay
%   D8.4 SASE/ZTNA explicitly NOT in the reference, with documented rationale
%   D8.5 Network sovereignty scores
% ============================================================================

\chapter{Network and segmentation}
\label{ch:network}

\begin{caveatbox}[Dependencies]
\noindent This chapter applies the layer template of
\trchap{ch:identity}{} to the network layer. It consumes policy
decisions from \trchap{ch:policy}{} (in particular for attachment
authorisation and east-west allowance), it produces flow records
consumed by \trchap{ch:observability}{}, and it relies on standards
from \S\ref{sec:ap-standards}. The layer belongs to the \emph{Execution plane}.
\end{caveatbox}

The network layer answers a question that is operationally simple
and architecturally consequential: \emph{which traffic is permitted
between which parties}. It is the first layer in which the
institution's sovereignty posture meets its physical infrastructure
--- cables, fibres, switches, wireless access points --- and the
layer in which the longest-lived investments tend to reside. A
sovereignty analysis that succeeds at the identity and policy planes
but fails at the network plane has accomplished little of value.

%-----------------------------------------------------------------------------
\section{Functional role and interface contract}
\label{sec:nw-contract}

The network layer is responsible for: declaring and operating
zones with documented security profiles; authenticating and
attaching subjects to the appropriate zone on the basis of identity
and posture; enforcing east-west and north-south traffic rules
returned by the policy layer; emitting flow telemetry of sufficient
fidelity for the observability layer to support audit and incident
response; and providing, where the institution's research mission
requires unobstructed high-throughput data movement, a research-data
path that the general stateful-filtering perimeter would otherwise
penalise --- realised, according to the institution's connectivity, as
a Science DMZ, a dedicated circuit or overlay operated by the
institution's \gls{nren}, or
a tuned Data-Transfer-Node fabric~\cite{geant-dtn}.

\begin{contractbox}[Network layer]
\noindent A compliant network layer, in the sense of this
architecture, must:

\begin{itemize}[leftmargin=1.5em,nosep]
  \item declare a zone catalogue with, for each zone, a published
    security profile (filtering rules, egress policy, audit
    obligations, posture requirements) and a \keyword{zone
    responsibility charter} --- the dossier's \emph{enclave
    charter} --- naming the accountable owner of that profile;
  \item authenticate and attach subjects via
    RADIUS~\cite{rfc-2865-radius} using attributes that carry, at
    minimum, the role activated and the assigned zone;
  \item configure devices through NETCONF~\cite{rfc-6241-netconf} /
    YANG~\cite{rfc-7950-yang} \emph{or} an equivalent declarative
    interface, without dependence on proprietary command-line
    interfaces in critical paths;
  \item exchange routing information through standard protocols
    (BGP, EVPN where applicable);
  \item emit flow records in IPFIX~\cite{rfc-7011-ipfix} format with
    sufficient fields for the observability layer to perform
    correlation;
  \item express east-west and north-south traffic rules in a form
    that can be reviewed, version-controlled, and reconciled against
    the decisions emitted by \trchap{ch:policy}{};
  \item provide (or document an architectural decision against) a
    high-throughput research-data path not penalised by the general
    stateful-filtering perimeter --- a \gls{sciencedmz}, a dedicated
    circuit or lightpath operated by the institution's \gls{nren},
    a dedicated overlay
    (research-community L3VPN) connection, or a tuned Data-Transfer-Node
    fabric~\cite{geant-dtn} --- kept identity- and posture-attached,
    IDS-monitored, policy-governed and flow-mirrored to observability
    (the stateful firewall is bypassed, not the trust model).
\end{itemize}
\end{contractbox}

The contract is silent on the physical topology, on the choice of
overlay technology, and on the specific zones declared --- which
depend on the institution's mix of administration, teaching,
research, and shared services. The contract requires that whatever
choices are made be documented, governed by the policy layer, and
visible to the observability layer.

%-----------------------------------------------------------------------------
\section{Reference instantiation: VLAN/VRF + FreeRADIUS + Open vSwitch + Science DMZ}
\label{sec:nw-reference}

The reference instantiation combines three patterns of differing
maturity: a classical VLAN / VRF physical segmentation; a FreeRADIUS
attachment service driven by the institution's identity and policy
layers; an Open vSwitch overlay for the cases where physical
segmentation alone is insufficient; and a Science DMZ for the
research data flows that the enterprise filtering perimeter would
otherwise penalise.

\begin{instantiationbox}[Network layer]
\noindent\textbf{Topology.} The campus fabric is partitioned at
layer~2 into \gls{vlan}s per zone (administration, teaching,
research, BYOD, IoT, sensitive, management, Science DMZ); routing
between zones is held in layer-3 \emph{virtual routing and
forwarding} (VRF) instances with explicit transit rules. Wired and
wireless attachment is authenticated by FreeRADIUS~\cite{freeradius}, which receives
attributes from the identity provider (\trchap{ch:identity}{}) and
queries the policy layer (\trchap{ch:policy}{}) to determine the
attachment zone and per-subject allowances. Where micro-segmentation
within a zone is required (for example, between research projects
sharing a single VLAN), Open vSwitch~\cite{openvswitch} flows are programmed by an SDN
controller from policy-layer outputs. The Science DMZ is a
dedicated subnet bypassing the enterprise filtering perimeter,
protected by IDS and access-control lists rather than by stateful
firewalling; flows are mirrored to the observability layer.
\end{instantiationbox}

\paragraph{Research-network-operated realisations.} The same function
is realised at national and global scale through research-network- and
experiment-operated dedicated paths rather than a campus-perimeter
enclave: the regional and national \glsshort{nren}s
provide dedicated circuits and lightpaths, and large research
collaborations operate dedicated optical private networks and L3VPN
overlays~\cite{cern-lhcone, cern-wlcg} --- a dedicated, segmented
research-data fabric connecting the originating facility and its
distributed data tiers. The performance-measurement component,
perfSONAR~\cite{perfsonar}, is co-maintained by the major
\glsshort{nren}s and is common to all of these realisations. The
contract of \S\ref{sec:nw-contract} is satisfied by any of them; the
Science DMZ is the open campus-scale reference, not the only form.

\paragraph{Zone catalogue.} The reference catalogue, adapted from
established practice in research-intensive institutions, includes
the zones in Table~\ref{tab:nw-zones}. The catalogue is illustrative;
each institution declares its own. The contractual obligation is
that the catalogue exist in writing and that each zone have a
documented security profile.

\begin{table}[H]
\centering\small
\caption{Reference zone catalogue (illustrative). Each row in an
institutional catalogue is owned by a named accountable team.}
\label{tab:nw-zones}
\begin{tabularx}{\textwidth}{%
  >{\hsize=0.5\hsize\linewidth=\hsize\bfseries\raggedright\arraybackslash}X%
  >{\hsize=1.5\hsize\linewidth=\hsize\raggedright\arraybackslash}X%
}
\toprule
Zone & Security profile (illustrative) \\
\midrule
Administration & Stateful filtering, MFA-required for access, restricted egress, full audit. \\
Teaching       & Stateful filtering, role-based access, broad egress, standard audit. \\
Research (general) & Stateful filtering, role-and-project based access, project-scoped egress, audit by project. \\
Research (high-throughput data path; e.g.\ Science DMZ) & ACL-based filtering only, IDS, dedicated IPv4/IPv6 prefixes, perfSONAR measurement, egress by IP prefix. \\
BYOD           & Captive portal with \gls{nac} posture verification, Internet egress only, isolated from internal zones. \\
IoT            & Network-only egress to documented destinations, no internal east-west, periodic posture re-attestation. \\
Sensitive      & Stricter than Administration; key custody and disclosure regime documented; no direct Internet egress. \\
Management     & Reachable only from administrative jump hosts; audited at every connection. \\
\bottomrule
\end{tabularx}
\end{table}

\paragraph{Operational considerations.} The reference instantiation
is the most demanding of the layer chapters in terms of skills: it
would typically require network engineering (VLAN, VRF, BGP), Linux
administration (FreeRADIUS, OvS), and a working command of the
policy layer (\trchap{ch:policy}{}). For a research-intensive
institution, a small team of engineers shared with adjacent duties
is an indicative envelope; actual figures vary with the campus
footprint and the research-network mission. The labour market is
deep at the level of
standard protocols and narrower at the level of integrating them
into a TAC-aware attachment service.

\begin{realworldbox}[Science DMZ adoption at scale]
\noindent The Science DMZ pattern was developed within the
high-performance research-networking community and is the
subject of detailed public documentation. The pattern has been
implemented at a number of research-intensive universities and at
large-facility experimental boundaries. Large university networks
would typically operate multi-zone architectures with explicit
security profiles per zone, using a combination of proprietary and
open components --- the brownfield hybrid pattern that this Reference's
playbook anticipates. Specific zoning details and current deployment
status of any named institution should be verified against current
public documentation before being cited.
\end{realworldbox}

\begin{realworldbox}[Experiment-operated optical private networks and research backbones]
\noindent At national and global scale the same need is met by
research-network- and experiment-operated dedicated networks rather
than by a campus enclave. A large-collaboration optical private
network (originating facility to the principal data tiers) and a
companion L3VPN overlay (any-to-any between the data
tiers)~\cite{cern-lhcone, cern-wlcg} carry experiment data on a
dedicated, segmented fabric; the regional and national
\glsshort{nren}s~\cite{geant} provide the backbone and
dedicated-circuit services on which member networks build, and
research-network testbed work on Data Transfer
Nodes~\cite{geant-dtn} addresses the same high-throughput transfer
problem at the host level. These realise the same contract function as
the Science DMZ, organised as a managed dedicated network rather than
a perimeter-bypass subnet.
\end{realworldbox}

%-----------------------------------------------------------------------------
\section{Alternative~1 --- Proprietary macro-segmentation speaking standards}
\label{sec:nw-alt1}

\begin{alternativebox}[Proprietary macro-segmentation]
\noindent The campus fabric is built on Cisco, Juniper, Aruba, or
Arista equipment, configured through standard protocols (RADIUS for
attachment, NETCONF/YANG for configuration, BGP/EVPN for routing).
The institution retains the contracts of \S\ref{sec:nw-contract}
while using established commercial gear.
\end{alternativebox}

\paragraph{Where it adds value.} The skills market is the deepest of
all network alternatives; the gear is mature and well supported;
integration with existing operational tooling (monitoring, asset
management, change management) is direct.

\paragraph{Where it constrains the institution.} Two recurring
constraints. First, the supplier's management consoles often expose
features that the standard interfaces do not: an institution that
adopts the management console for non-trivial operations couples
itself to the supplier even though the underlying configuration is
notionally NETCONF-based. Second, premium features (behavioural
analytics, advanced segmentation, integrated identity) are licensed
separately and shift between licence tiers across renewals; their
adoption in critical paths erodes replaceability.

\paragraph{When it is the right choice.} Most institutions, in
practice, will operate the network on commercial gear for at least
the medium term. The architectural commitment is not to refuse
proprietary equipment but to keep it on the standard side of its
interfaces, treat its management consoles as advisory, and use only
features that the institution can also obtain from at least one
other supplier.

%-----------------------------------------------------------------------------
\section{Alternative~2 --- Micro-segmentation overlay}
\label{sec:nw-alt2}

\begin{alternativebox}[Micro-segmentation overlay]
\noindent Segmentation is implemented as a software overlay
independent of the underlying physical fabric: VMware NSX, Illumio,
Cisco ACI, or Cilium ClusterMesh in Kubernetes-heavy environments.
Policy is expressed at the workload level, often per-process or
per-container, and enforced at host or hypervisor.
\end{alternativebox}

\paragraph{Where it adds value.} Granularity. Micro-segmentation
allows differentiation of traffic between workloads sharing a
physical zone; in environments where a research group's workloads
must be isolated from each other without provisioning per-group
VLANs, this is the right tool. Cilium-based~\cite{cilium} micro-segmentation in
Kubernetes is also the architectural fit for the compute layer
(\trchap{ch:compute}{}) when policy needs to traverse the K8s and
VM boundaries.

\paragraph{Where it constrains the institution.} Commercial
micro-segmentation overlays (NSX, Illumio, ACI) have substantial
licence costs and tend to express policy in supplier-specific
languages whose translation to other products is non-trivial; the
exit-rehearsal condition of \S\ref{sec:ap-replaceability} is most challenged here. Cilium
and other open alternatives are more recent and less broadly skilled
than their commercial counterparts.

\paragraph{When it is the right choice.} Institutions with a high
density of multi-tenant workloads (research clouds, multi-faculty
shared platforms) where macro-segmentation does not provide
sufficient granularity. The recommendation is to adopt the open
alternative (Cilium) as the reference, and the commercial
alternatives only when the institution's posture cost-justifies them.

%-----------------------------------------------------------------------------
\section{Discussion --- SASE and ZTNA: why they are not the reference}
\label{sec:nw-sase}

Secure Access Service Edge (SASE) and Zero Trust Network Access
(ZTNA) services from suppliers such as Cloudflare, Zscaler, and
Netskope present an attractive operational profile: policy
expressed centrally, enforcement near the user, deep integration
with identity, and reduced reliance on on-premise filtering
infrastructure. In a sovereignty-aware analysis, however, they
present a structural difficulty.

The data plane of a SASE service is operated by the supplier, in
locations and under jurisdictions chosen by the supplier, with
inspection capabilities that the supplier necessarily exercises (TLS
termination, content inspection, identity correlation). The supplier
becomes a recipient of the institution's network traffic in a sense
that the standards-only configurations of \S\S\ref{sec:nw-reference}
and \ref{sec:nw-alt1} do not. The data dimension of the sovereignty
grid is consequently bounded: a SASE deployment, in its standard
configuration, cannot exceed Partial on the disclosure-regime
indicator.

This is not an objection to SASE as a technology. SASE is an
appropriate choice in two bounded scenarios: as a complement to the
reference architecture for a specific population (for example,
unmanaged contractor devices) where the institution accepts the
sovereignty cost in exchange for a manageability gain; and as a
transitional measure during a brownfield migration whose end-state
is the reference architecture. The architectural posture is to
treat SASE as a tactical tool, not as the structural network layer.

\paragraph{Zero-Trust anchoring.} Declining SASE/ZTNA does not
forgo Zero Trust: the reference architecture realises the seven
tenets of NIST SP~800-207~\cite{nist-zt} structurally --- every
layer treated as a protected resource with communication secured
regardless of network location (Tenets~1--2); identity-mediated,
per-session, strictly-enforced access through the identity and
policy planes (Tenets~3 and~6); access decided by dynamic policy on
identity, role, posture and environment, with default-deny zones
(Tenet~4); continuous posture and integrity monitoring through the
endpoint and observability layers (Tenet~5); and decision/access
telemetry collected through the canonical TAC event schema to refine
policy (Tenet~7). Zero Trust is thus a property of the architecture,
not a procured service.

%-----------------------------------------------------------------------------
\section{Sovereignty grid applied}
\label{sec:nw-sovereignty}

\begin{table}[H]
\centering\small
\caption{Per-dimension sovereignty profile of the network
instantiations.}
\label{tab:nw-sovereignty}
\begin{tabularx}{\textwidth}{%
  >{\hsize=0.6\hsize\linewidth=\hsize\bfseries\raggedright\arraybackslash}X%
  >{\hsize=1.1\hsize\linewidth=\hsize\centering\arraybackslash}X%
  >{\hsize=1.3\hsize\linewidth=\hsize\centering\arraybackslash}X%
  >{\hsize=1.4\hsize\linewidth=\hsize\centering\arraybackslash}X%
  >{\hsize=0.6\hsize\linewidth=\hsize\centering\arraybackslash}X%
}
\toprule
Dimension & Reference (open) & Proprietary macro-seg & Micro-seg (commercial) & SASE \\
\midrule
Data          & 3 Robust      & 3 Robust      & 2 Established & 1 Partial \\
Technical     & 3 Robust      & 2 Established & 1 Partial     & 1 Partial \\
Financial     & 2 Established & 2 Established & 1 Partial     & 1 Partial \\
Operational   & 2 Established & 3 Robust      & 2 Established & 3 Robust  \\
Institutional & 3 Robust      & 2 Established & 1 Partial     & 1 Partial \\
\bottomrule
\end{tabularx}
\end{table}

\paragraph{Driving indicators by dimension.} Data: disclosure regime
and storage location. Technical: source visibility and presence of
captive extensions in the configuration path. Financial: cost
predictability and exit-cost proportionality. Operational: depth of
the skills market and quality of public documentation.
Institutional: roadmap influence and integration freedom.

\noindent The network layer illustrates clearly the trade-off the
sovereignty grid is designed to surface: each row of the table is a
different way of distributing capacity along the five dimensions.
The reference instantiation is dominant on Data, Technical and
Institutional; the proprietary macro-segmentation gains on
Operational at modest cost on Technical; the commercial
micro-segmentation and SASE deliver Operational and convenience but
at structural cost on Data and Institutional.

%-----------------------------------------------------------------------------
\section{Regulatory anchoring}
\label{sec:nw-regulatory}

\noindent The network layer is the perimeter and segmentation
substrate underpinning \gls{gdpr}~Art.~32(1) (security of
processing through technical measures). It operationalises
NIS2~Art.~21(2)(a) (risk analysis applied to network design),
(g) (cyber hygiene baseline including segregation), and indirectly
(b) (incident handling, through the chokepoints that bound
incident scope). The layer is the principal anchor of
27001~A.8.20 (network security) and A.8.22 (segregation of
networks); it contributes to NIST CSF \textsf{PROTECT} (PR.IR ---
infrastructure resilience) and \textsf{DETECT} (DE.CM, through
flow telemetry feeding the audit pipeline). Detailed mapping in
\trchap{ch:regulatory-mapping}{} \S\ref{sec:rm-gdpr},
\S\ref{sec:rm-nis2}, \S\ref{sec:rm-iso}, \S\ref{sec:rm-nist-csf}.

%-----------------------------------------------------------------------------
\section{Common pitfalls}
\label{sec:nw-pitfalls}

\begin{pitfallbox}[Standard interfaces promised, proprietary management used]
\noindent An institution adopts NETCONF/YANG-capable gear, then
operates it through the supplier's management console for
day-to-day work. Within two to three years, the operational
practices, the runbooks, and the team's reflexes are tied to the
console. The standards remain technically available but
operationally unused. The mitigation is to make the standard
interface the operational default, treat the management console as a
diagnostic aid only, and exercise the standard-only operation
periodically as part of the conformance suite of
\trchap{ch:conformance}{}.
\end{pitfallbox}

\begin{pitfallbox}[Science DMZ behind the enterprise firewall]
\noindent A Science DMZ that traverses the enterprise stateful
filtering perimeter loses its purpose: the throughput penalty that
the Science DMZ exists to avoid is reintroduced. The mitigation is
explicit: the Science DMZ has its own access path to the Internet,
protected by ACLs and IDS but not by a stateful firewall, and the
architectural decision to bypass the enterprise perimeter is
documented and reviewed.
\end{pitfallbox}

\begin{pitfallbox}[East-west allow-by-default]
\noindent Many campus networks permit east-west traffic by default
within a VLAN; consequently, the segmentation is at the zone
boundary only. Lateral movement during an incident is then bounded
only by the zone, not by the workload. The mitigation is to treat
intra-zone east-west as a policy decision, expressed in
\trchap{ch:policy}{} and enforced at the appropriate layer (host
firewall, micro-segmentation overlay, or hypervisor); the default is
deny, with allow-listed flows.
\end{pitfallbox}

\begin{pitfallbox}[ACL drift]
\noindent Access control lists are added when needed and removed
rarely. Within a few years, the ACL set on each device is a
historical record of past needs rather than a current expression of
policy. The mitigation is to express ACLs as the materialisation of
policy decisions (\trchap{ch:policy}{}) under version control, with
periodic reconciliation against the policy corpus and removal of
rules whose justification has lapsed.
\end{pitfallbox}

\begin{pitfallbox}[BYOD VLAN as ``isolation'' without enforcement]
\noindent A separate BYOD VLAN is declared but the enforcement that
prevents BYOD endpoints from reaching internal resources is missing
or broken. Users discover this when they encounter problems and
their resourceful colleagues find paths around the intended
isolation. The mitigation is to make the BYOD posture an enforceable
policy, verified by NAC, with automated periodic re-attestation, and
to monitor the actual flows from the BYOD VLAN against the policy.
\end{pitfallbox}
